\section{Software Architecture}
\label{sec:software_architecture}

The firmware is built upon a non-blocking, event-driven architecture.
This approach avoids busy-waiting and blocking delays, ensuring responsive \ac{MCU} performance even when handling complex \ac{LED} animations, rotary encoder inputs, or network communications.

\subsection{Dual-Core Task Partitioning}
\label{sec:dual_core}

The ESP32-S3 integrates two Xtensa~LX7 processing cores.
The firmware exploits this by statically partitioning work at initialisation: each core runs a dedicated FreeRTOS task pinned via \texttt{xTask\-Create\-Pinned\-To\-Core()}, so the real-time peripheral path and the latency-tolerant user interface never compete for the same execution context.

\begin{figure}[H]
  \centering
  \begin{circuitikz}[>=latex]
    % Core 0 block
    \draw[dashed, thick, fill=FrostBg, rounded corners] (0.0, 0.0) rectangle (5.4, 5.4);
    \node[anchor=north, font=\bfseries] at (2.7, 5.4) {Core 0 — Real-Time};
    \node[draw, fill=white, minimum width=4.4cm, minimum height=0.7cm, rounded corners, font=\small, align=center] at (2.7, 4.5) {WLAN Stack};
    \node[draw, fill=white, minimum width=4.4cm, minimum height=0.7cm, rounded corners, font=\small, align=center] at (2.7, 3.4) {Art-Net Parser};
    \node[draw, fill=white, minimum width=4.4cm, minimum height=0.7cm, rounded corners, font=\small, align=center] at (2.7, 2.3) {DMX-512 Receiver};
    \node[draw, fill=white, minimum width=4.4cm, minimum height=0.7cm, rounded corners, font=\small, align=center] at (2.7, 1.2) {LED Controller (RMT)};

    % Shared config block
    \node[draw, fill=white, minimum width=3.0cm, minimum height=1.8cm, align=center, rounded corners] (cfg) at (8.0, 2.7) {\small\texttt{lighting\_config\_t}\\\small\texttt{portMUX\_TYPE}};

    % Core 1 block
    \draw[dashed, thick, fill=FrostBg, rounded corners] (10.8, 0.0) rectangle (16.2, 5.4);
    \node[anchor=north, font=\bfseries] at (13.5, 5.4) {Core 1 — User Interface};
    \node[draw, fill=white, minimum width=4.4cm, minimum height=0.7cm, rounded corners, font=\small, align=center] at (13.5, 4.5) {Event Manager};
    \node[draw, fill=white, minimum width=4.4cm, minimum height=0.7cm, rounded corners, font=\small, align=center] at (13.5, 3.6) {Timer Manager};
    \node[draw, fill=white, minimum width=4.4cm, minimum height=0.7cm, rounded corners, font=\small, align=center] at (13.5, 2.7) {Rotary Encoder / Button};
    \node[draw, fill=white, minimum width=4.4cm, minimum height=0.7cm, rounded corners, font=\small, align=center] at (13.5, 1.8) {OLED Display (I\textsuperscript{2}C)};
    \node[draw, fill=white, minimum width=4.4cm, minimum height=0.7cm, rounded corners, font=\small, align=center] at (13.5, 0.9) {NVS Flash Storage};

    % Bidirectional arrows through shared config
    \draw[latex-latex, thick] (5.4, 2.7) -- (cfg.west);
    \draw[latex-latex, thick] (cfg.east) -- (10.8, 2.7);
  \end{circuitikz}
  \caption{Static task partitioning across the two ESP32-S3 cores.}
  \label{fig:dual_core_arch}
\end{figure}

\paragraph{Core 0 — Real-Time Peripherals}
Core~0 hosts all tasks that must execute with minimal and predictable latency.
The ESP-IDF Wi-Fi driver is bound to Core~0 by default, making it the natural home for the entire network-facing path.
Art-Net \ac{UDP} packets are received and parsed directly on this core.
The \ac{DMX}-512 receiver services the incoming \ac{UART} stream and writes the working frame buffer without blocking.
Computed pixel data is driven to the \ac{LED} string via the \ac{RMT} peripheral, also from Core~0, to guarantee consistent frame timing.
No blocking calls or indefinite waits are permitted on this core.

\paragraph{Core 1 — User Interface}
Core~1 handles all tasks that tolerate occasional scheduling jitter.
The Event Manager and Timer Manager (described in the following subsections) both execute here, along with debounced rotary encoder and button handling.
\ac{OLED} display updates are issued over \ac{I2C} exclusively from Core~1, so the variable latency of \ac{I2C} transactions never affects the real-time path.
Non-volatile parameter reads and writes to the ESP-IDF \ac{NVS} flash \ac{API} are similarly confined to this core.

\subsubsection{Shared Configuration and Spinlock Protection}
\label{sec:shared_config}

Both cores must read and occasionally write a small set of operating parameters: the current operating mode, \ac{DMX} start address, and Art-Net universe.
These are consolidated into a single \texttt{lighting\_config\_t} struct, keeping the shared data surface minimal and explicit.

Calling \texttt{portENTER\_CRITICAL()} without an argument only masks interrupts on the \emph{calling} core and provides no protection against simultaneous access from the other core.
The correct dual-core primitive is \texttt{portMUX\_TYPE}, which wraps the ESP32-S3 hardware spinlock.
When one core acquires the lock, the other core busy-waits until it is released.
Both cores must therefore hold the lock for the shortest possible time — in this implementation, strictly to a struct copy with no processing inside the critical section.

\begin{longlisting}
\begin{minted}{cpp}
// Shared declaration (single translation unit)
portMUX_TYPE config_mux = portMUX_INITIALIZER_UNLOCKED;
lighting_config_t lighting_config;	// protected by config_mux

// Core 1 — writing a new configuration
void config_set(const lighting_config_t* src) {
  portENTER_CRITICAL(&config_mux);
  lighting_config = *src;	// struct copy only, no processing
  portEXIT_CRITICAL(&config_mux);
}

// Core 0 — reading a snapshot before processing a frame
lighting_config_t config_get(void) {
  lighting_config_t snap;
  portENTER_CRITICAL(&config_mux);
  snap = lighting_config; // struct copy only, no processing
  portEXIT_CRITICAL(&config_mux);
  return snap;
}
\end{minted}
\caption{Dual-core safe access to \texttt{lighting\_config\_t} using \texttt{portMUX\_TYPE}.}
\label{lst:spinlock_config}
\end{longlisting}

\begin{tcolorbox}[colback=gray!8, colframe=gray!50, boxrule=0.4mm, arc=5px, title=\textbf{Planned Enhancement — Queue-Based Configuration Updates}]
As firmware complexity grows, a more robust pattern is for Core~1 to \emph{own} the authoritative copy of \texttt{lighting\_config\_t} and post change requests to Core~0 via a minimal single-element FreeRTOS queue.
Core~0 applies the update at a safe point between rendered frames, removing the spinlock from the real-time path entirely and eliminating the theoretical stall on Core~0 caused by a Core~1 write occurring mid-frame.
The spinlock approach is retained for now because configuration changes are infrequent and the critical section is provably restricted to a single struct copy.
\end{tcolorbox}

\subsection{Event Manager}
\label{sec:event_manager}

The Event Manager decouples event generation from event processing on Core~1.
Event sources — the rotary encoder \ac{ISR}, the button debounce handler, and the Timer Manager — push events into a lightweight array-based ring buffer.
The Core~1 main loop polls this buffer and passes each event to the \texttt{process\_event()} handler, which drives the menu \ac{FSM}.

Each event consists of an \texttt{event\_id\_t} identifier and an optional \SI{32}{\bit} parameter.
The ring buffer capacity is sufficient for the expected maximum burst rate of user inputs and timer expirations; no priority discrimination is required at the current firmware complexity level, since time-critical work (network reception, \ac{LED} output) is handled entirely on Core~0.

\begin{figure}[H]
  \centering
  \begin{circuitikz}[>=latex]
    % Event sources
    \node[draw, fill=white, minimum width=2.8cm, minimum height=0.8cm, rounded corners] (enc) at (0, 3.0) {Rotary Encoder ISR};
    \node[draw, fill=white, minimum width=2.8cm, minimum height=0.8cm, rounded corners] (btn) at (0, 1.8) {Button Handler};
    \node[draw, fill=white, minimum width=2.8cm, minimum height=0.8cm, rounded corners] (tim) at (0, 0.6) {Timer Manager};

    % Ring buffer
    \node[draw, fill=white, minimum width=2.8cm, minimum height=2.6cm, align=center, rounded corners] (rb) at (5.0, 1.8) {Ring Buffer\\\small(array-based FIFO)};

    % Menu engine
    \node[draw, fill=white, minimum width=2.8cm, minimum height=0.8cm, rounded corners] (menu) at (10.0, 1.8) {Menu \ac{FSM}};

    % Arrows
    \draw[-latex, thick, rounded corners] (enc.east) -- ++(0.6, 0) |- (rb.west);
    \draw[-latex, thick] (btn.east) -- (rb.west);
    \draw[-latex, thick, rounded corners] (tim.east) -- ++(0.6, 0) |- (rb.west);
    \draw[-latex, thick] (rb.east) -- (menu.west);
  \end{circuitikz}
  \caption{Event Manager data flow on Core~1.}
  \label{fig:event_manager_arch}
\end{figure}

The \texttt{process\_event()} handler evaluates the current menu depth and state, then routes the event to the appropriate response: an \ac{OLED} display refresh, a \ac{NVS} parameter write, or a \texttt{lighting\_config\_t} update via the spinlock accessor described in Section~\ref{sec:shared_config}.

\begin{longlisting}
\begin{minted}{cpp}
void process_event(const event_t* event) {
  uint8_t depth = main_sw.get_deepness();
  uint8_t state = main_sw.get_current();

  switch (event->id) {
    case EVT_ENCODER_STEP:
      menu_navigate(event->param);   // signed step value
      break;
    case EVT_BUTTON_CLICK:
      menu_select(depth, state);
      break;
    case EVT_DISPLAY_REFRESH:
      oled_render();
      break;
    // Additional event handling...
  }
}
\end{minted}
\caption{Event dispatch into the menu \ac{FSM}.}
\label{lst:process_event}
\end{longlisting}

\subsection{Software Timer Manager}
\label{sec:timer_manager}

On Core~1, the \texttt{SoftTimerManager} class handles all non-blocking timing requirements.
It stores a fixed-size static array of timer entries and scans them on every iteration of the Core~1 main loop, comparing each entry's deadline against the current \texttt{millis()} value.
No hardware timer or \ac{RTOS} timer resource is consumed; the \texttt{millis()} counter is the sole time reference.

When a timer expires, it pushes a predefined \texttt{event\_id\_t} into the Event Manager ring buffer.
The Core~1 event loop then picks up the event and forwards it to \texttt{process\_event()} in the same way as any input event.
For example, the display standby timer pushes \texttt{EVT\_DISPLAY\_STANDBY} when it expires, and the menu \ac{FSM} reacts by blanking the \ac{OLED}.

\begin{figure}[H]
  \centering
  \begin{circuitikz}[>=latex]
    % Timer Manager block
    \draw[dashed, thick, fill=FrostBg, rounded corners] (2.8, 0.5) rectangle (8.6, 3.5);
    \node[anchor=north] at (5.7, 3.5) {\textbf{Timer Manager}};

    % Input
    \node[draw, fill=white, minimum width=2.8cm, minimum height=0.8cm, rounded corners] (loop) at (0, 2.0) {Core~1 Main Loop};

    % Internal
    \node[draw, fill=white, minimum width=3.2cm, minimum height=0.8cm, rounded corners] (arr) at (5.7, 2.6) {Static Timer Array};
    \node[draw, fill=white, minimum width=3.2cm, minimum height=0.8cm, rounded corners] (cmp) at (5.7, 1.4) {Expiry Check (\texttt{millis})};

    % Output
    \node[draw, fill=white, minimum width=2.8cm, minimum height=0.8cm, rounded corners] (rb) at (11.2, 2.0) {Event Ring Buffer};

    % Arrows
    \draw[-latex, thick] (loop.east) -- (2.8, 2.0);
    \draw[-latex, thick] (arr.south) -- (cmp.north);
    \draw[-latex, thick, rounded corners] (cmp.east) -- ++(0.6, 0) |- (8.7, 2.0) -- (rb.west);
  \end{circuitikz}
  \caption{Software Timer Manager data flow — static array polled each Core~1 loop iteration.}
  \label{fig:timer_manager_arch}
\end{figure}

\begin{table}[H]
  \caption{Software timer operational modes.}
  \label{tab:timer_modes}
  \centering
  \begin{tabularx}{\textwidth}{|l|X|}
    \hline
    \textbf{Mode} & \textbf{Description} \\ \hline \hline
    Periodic & Timer restarts automatically after posting its event. Used for animations and cyclic updates. \\ \hline
    One-shot & Timer halts after expiration. Used for timeouts, such as the display standby or fan cooldown. \\ \hline
  \end{tabularx}
\end{table}


\subsection{Menu Engine}
\label{sec:menu_engine}

The menu engine implements a data-driven, tree-based user interface on Core~1 with zero dynamic memory allocation.
All menu structure is encoded in \texttt{const} arrays of structs that the linker places directly in the ESP32-S3's flash memory; no heap allocations are performed at runtime, and no \ac{RAM} is consumed beyond a single pointer to the currently active page.

\subsubsection{Static Data Model}
\label{sec:menu_data_model}

The engine operates on two mutually referential types: \texttt{menu\_page\_t} and \texttt{menu\_item\_t}.
A page represents one visible layer of the menu hierarchy; an item represents one selectable entry within a page.
Each item carries a text label, a pointer to a flash-resident icon bitmap, a type discriminator, and a union target interpreted by the engine according to the discriminator.

\begin{table}[H]
  \caption{Menu item types and their engine behaviour.}
  \label{tab:menu_item_types}
  \centering
  \begin{tabularx}{\textwidth}{|l|X|}
    \hline
    \textbf{Type} & \textbf{Engine Behaviour} \\ \hline \hline
    \texttt{ITEM\_SUBMENU}     & Navigates to the page pointed to by \texttt{target.submenu}. \\ \hline
    \texttt{ITEM\_ACTION}      & Calls the void function referenced by \texttt{target.action\_cb}. \\ \hline
    \texttt{ITEM\_EDIT\_UINT16}& Enters an in-place value editor bounded by \texttt{min\_val} and \texttt{max\_val}; the encoder wraps the value modularly within the defined range. \\ \hline
    \texttt{ITEM\_BACK}        & Returns to the page referenced by the current page's \texttt{parent} pointer. \\ \hline
  \end{tabularx}
\end{table}

Editable numeric parameters use the \texttt{edit\_uint16\_t} sub-struct, which holds a direct pointer to the target global variable alongside its permitted bounds.
Incrementing past \texttt{max\_val} wraps the value back to \texttt{min\_val}, and decrementing below \texttt{min\_val} wraps forward to \texttt{max\_val}.

\begin{longlisting}
\begin{minted}{c}
// Forward declarations — required because menu_item_t and menu_page_t reference each other
typedef struct menu_page_t menu_page_t;
typedef struct menu_item_t menu_item_t;

// Icon wrapper — stores bitmap dimensions, rendering offset, and flash data pointer
typedef struct {
  uint8_t width;
  uint8_t height;
  int8_t  x_offset;
  int8_t  y_offset;
  const uint8_t* data;
} icon_t;

// Item type discriminator
typedef enum {
  ITEM_SUBMENU,      // Opens a deeper menu page
  ITEM_ACTION,       // Executes a void callback
  ITEM_EDIT_UINT16,  // In-place numeric editor with bounds
  ITEM_BACK          // Returns to the parent page
} menu_item_type_t;

// Bounds descriptor for editable uint16 parameters
typedef struct {
  uint16_t* val_ptr;   // Pointer to the variable being edited
  uint16_t  min_val;
  uint16_t  max_val;
} edit_uint16_t;
\end{minted}
\caption{Supporting type definitions: icon wrapper, item type discriminator, and bounds descriptor.}
\label{lst:menu_types_supporting}
\end{longlisting}

\begin{longlisting}
\begin{minted}{c}
// Menu item — one selectable entry within a page
struct menu_item_t {
  const char*      label;
  const icon_t*    icon;
  menu_item_type_t type;
  union {
    menu_page_t* submenu;
    void        (*action_cb)(void);
    edit_uint16_t edit;
  } target;
};

// Menu page — one visible layer of the hierarchy
struct menu_page_t {
  const char*  title;
  menu_page_t* parent;    // NULL for the top-level page
  uint8_t      num_items;
  menu_item_t* items;
};
\end{minted}
\caption{Composite types: \texttt{menu\_item\_t} and \texttt{menu\_page\_t}.}
\label{lst:menu_types_composite}
\end{longlisting}

\subsubsection{Menu Tree Instantiation}
\label{sec:menu_tree}

Menu trees are built bottom-up in C: leaf pages must be defined before the parent pages that reference them.
The top-level page carries a \texttt{NULL} parent pointer; all other pages carry a pointer to their direct ancestor, enabling \texttt{ITEM\_BACK} navigation without maintaining a runtime stack.

\begin{figure}[H]
  \centering
  \begin{circuitikz}[>=latex]
    % Main page node
    \node[draw, fill=white, minimum width=3.0cm, minimum height=0.8cm, rounded corners] (main) at (1.5, 2.0) {\texttt{page\_main}};

    % Layer-2 nodes
    \node[draw, fill=white, minimum width=3.2cm, minimum height=0.8cm, rounded corners] (artnet) at (7.5, 3.2) {\texttt{page\_artnet}};
    \node[draw, fill=white, minimum width=3.2cm, minimum height=0.8cm, rounded corners] (dmxp)   at (7.5, 2.0) {\texttt{page\_dmx}};
    \node[draw, dashed, fill=FrostBg, minimum width=3.2cm, minimum height=0.8cm, rounded corners] (reb) at (7.5, 0.8) {\small\texttt{action\_reboot\_device}};

    % Navigation arrows
    \draw[-latex, thick, rounded corners] (main.east) -- ++(0.5,0) |- (artnet.west)
      node[near end, above, font=\scriptsize] {SUBMENU};
    \draw[-latex, thick] (main.east) -- (dmxp.west)
      node[midway, above, font=\scriptsize] {SUBMENU};
    \draw[-latex, thick, rounded corners] (main.east) -- ++(0.5,0) |- (reb.west)
      node[near end, below, font=\scriptsize] {ACTION};
  \end{circuitikz}
  \caption{Two-level menu tree for the RLCV4 firmware.}
  \label{fig:menu_tree}
\end{figure}

\begin{longlisting}
\begin{minted}{c}
// Global parameters referenced by the edit items
volatile uint16_t g_dmx_address = 1;
volatile uint16_t g_artnet_universe = 0;

// Action callbacks
void action_save_to_nvs(void)   { /* write to NVS */ }
void action_reboot_device(void) { /* esp_restart() */ }

extern menu_page_t page_main;

// ---- Layer 2: Art-Net submenu ----
menu_item_t items_artnet[] = {
  { "Universe", &icon_artnet, ITEM_EDIT_UINT16,
    .target.edit = { (uint16_t*)&g_artnet_universe, 0, 32767 } },
  { "Save NVS", &icon_save,   ITEM_ACTION,
    .target.action_cb = action_save_to_nvs },
  { "Back",     &icon_back,   ITEM_BACK,
    .target = {0} }
};
menu_page_t page_artnet = { "Art-Net Setup", &page_main, 3, items_artnet };

// ---- Layer 2: DMX submenu ----
menu_item_t items_dmx[] = {
  { "Address",  &icon_dmx,  ITEM_EDIT_UINT16,
    .target.edit = { (uint16_t*)&g_dmx_address, 1, 512 } },
  { "Save NVS", &icon_save, ITEM_ACTION,
    .target.action_cb = action_save_to_nvs },
  { "Back",     &icon_back, ITEM_BACK,
    .target = {0} }
};
menu_page_t page_dmx = { "DMX Setup", &page_main, 3, items_dmx };

// ---- Layer 1: Main menu ----
menu_item_t items_main[] = {
  { "Art-Net", &icon_artnet, ITEM_SUBMENU,
    .target.submenu = &page_artnet },
  { "DMX512",  &icon_dmx,   ITEM_SUBMENU,
    .target.submenu = &page_dmx },
  { "Reboot",  &icon_save,  ITEM_ACTION,
    .target.action_cb = action_reboot_device }
};
menu_page_t page_main = { "Main Menu", NULL, 3, items_main };
\end{minted}
\caption{Bottom-up instantiation of a two-level menu tree in flash memory.}
\label{lst:menu_tree}
\end{longlisting}

\subsection{Non-Volatile Parameter Storage}
\label{sec:nvs_storage}

To persist configuration parameters across power cycles, the firmware utilizes the ESP-IDF's native Non-Volatile Storage (\ac{NVS}) \ac{API}.
This subsystem operates on a dedicated partition in the internal flash memory and provides a wear-leveling, key-value storage mechanism that is robust against unexpected power loss during write operations.
This approach is superior to direct \ac{EEPROM} emulation, which lacks these safeguards.
All parameter persistence is handled on Core~1 to avoid impacting real-time tasks. \\

The save operation is triggered exclusively by a double-click of the rotary encoder button while navigating within a submenu.
When triggered, the firmware gathers all configurable runtime states into a single structure.
An 8-bit \ac{CRC} is then calculated across this data structure.
Finally, the structure and its \ac{CRC} are committed to flash as a single binary object. \\

On system boot, the firmware performs the reverse process to load the configuration.
It attempts to read the configuration from \ac{NVS}.
If it's found, the firmware recalculates the \ac{CRC} on the retrieved data portion and compares it to the stored \ac{CRC} value. \\

If the checksums match, the parameters are considered valid and are loaded into the active runtime configuration.
If the \ac{NVS} configuration is not found, is corrupted, or the \ac{CRC} check fails, the firmware discards the invalid data and loads a factory default configuration.
This ensures the device always starts in a known-good, operational state.

\begin{longlisting}
\begin{minted}{c}
// Container for storing configuration data and its integrity check
struct nvs_payload_t {
  config_t config_data; // All runtime parameters
  uint8_t  crc;         // CRC calculated over config_data
};

// Simplified logic for loading configuration from NVS on boot
void load_config_from_nvs(void) {
  nvs_payload_t payload;

  // Attempt to read the entire payload blob from NVS.
  // nvs_read_block() is a placeholder for the actual NVS API call.
  if (nvs_read_block(&payload, sizeof(payload))) {

    // config was read, now validate it.
    uint8_t calculated_crc = crc8_compute(
      (uint8_t*)&payload.config_data,
      sizeof(payload.config_data)
    );

    if (calculated_crc == payload.crc) {
      // CRC matches, apply the loaded configuration.
      apply_runtime_config(payload.config_data);
      return; // Success
    }
  }

  // If read failed, blob was wrong size, or CRC was invalid,
  // apply factory defaults to ensure a safe state.
  apply_factory_defaults();
}
\end{minted}
\caption{Simplified logic for reading and validating parameters from \ac{NVS} on boot.}
\label{lst:nvs_read_logic}
\end{longlisting}

\subsection{Network Lighting Protocol Receiver}
\label{sec:net_receiver}

Both Art-Net and \ac{sACN} (ANSI~E1.31) transport \ac{DMX}-512 channel data over standard \ac{IP} networks as \ac{UDP} datagrams.
Despite differing in packet format, universe numbering, and connection-management semantics, both protocols deliver an identical payload: an ordered array of up to \num{512} 8-bit channel values assigned to a numbered universe.
The receiver subsystem is designed around a protocol-agnostic shared data layer that exploits this common ground.
Either protocol driver writes incoming channel data into the same double-buffered universe store; the lighting engine consumes that data without any knowledge of which protocol produced it.

\begin{figure}[H]
  \centering
  \begin{circuitikz}[>=latex]
    % Protocol Receivers
    \draw[dashed, thick, fill=FrostBg, rounded corners] (0.0, 2.6) rectangle (5.2, 5.8);
    \node[anchor=north, font=\bfseries\small] at (2.6, 5.8) {Protocol Receivers};
    \node[draw, fill=white, minimum width=4.0cm, minimum height=0.8cm, rounded corners, font=\small] (artnet) at (2.6, 4.7) {Art-Net Receiver};
    \node[draw, fill=white, minimum width=4.0cm, minimum height=0.8cm, rounded corners, font=\small] (sacn)   at (2.6, 3.5) {sACN Receiver};
    % Network Manager below the box
    \node[draw, fill=white, minimum width=4.0cm, minimum height=0.8cm, rounded corners, font=\small] (netmgr) at (2.6, 1.3) {Network Manager};
    % Shared Data Layer
    \draw[dashed, thick, fill=FrostBg, rounded corners] (5.9, 0.4) rectangle (11.1, 5.8);
    \node[anchor=north, font=\bfseries\small] at (8.5, 5.8) {Shared Data Layer};
    \node[draw, fill=white, minimum width=4.0cm, minimum height=0.8cm, rounded corners, font=\small] (buf0) at (8.5, 4.7) {Primary Frame Buffer};
    \node[draw, fill=white, minimum width=4.0cm, minimum height=0.8cm, rounded corners, font=\small] (buf1) at (8.5, 3.5) {Secondary Frame Buffer};
    \node[draw, fill=white, minimum width=4.0cm, minimum height=0.8cm, rounded corners, font=\small] (cfg)  at (8.5, 1.3) {\texttt{universe\_config\_t}};
    % Output
    \node[draw, fill=white, minimum width=3.6cm, minimum height=0.9cm, rounded corners, font=\small] (light) at (14.3, 4.3) {Lighting Engine};
    \node[draw, fill=white, minimum width=3.6cm, minimum height=0.8cm, rounded corners, font=\small] (led)   at (14.3, 2.7) {LED Output};
    % Data: receivers -> buffers (straight horizontal, both sides at same y)
    \draw[-latex, thick] (artnet.east) -- (buf0.west);
    \draw[-latex, thick] (sacn.east)   -- (buf1.west);
    % Data: buffers -> lighting engine
    \draw[-latex, thick, rounded corners=4pt] (buf0.east) -- ++(0.25,0) |- (light.west);
    \draw[-latex, thick, rounded corners=4pt] (buf1.east) -- ++(0.2,0) |- (light.west);
    % Config -> lighting engine
    \draw[-latex, thick, rounded corners=4pt] (cfg.east) -- ++(1.3,0) |- (light.west);
    % Lighting -> LED
    \draw[-latex, thick] (light.south) -- (led.north);
    % Network Manager -> Protocol Receivers box (event control, dashed)
    \draw[-latex, dashed, thick] (netmgr.north) -- (2.6, 2.6);
  \end{circuitikz}
  \caption{Network lighting protocol receiver architecture.}
  \label{fig:net_receiver_arch}
\end{figure}

As shown in Figure~\ref{fig:net_receiver_arch}, solid arrows show the data path; the dashed arrow indicates event-based connection-state signalling from the Network Manager to the protocol receivers.
Both receivers share the same frame buffers and configuration structure; the lighting engine is fully decoupled from the active protocol.

\subsubsection{Protocol Overview}
\label{sec:net_proto_overview}

Table~\ref{tab:proto_comparison} summarises the operational properties of the two supported protocols.
Both deliver \num{512} \ac{DMX} channel values per universe as the innermost payload, making the shared data layer straightforward to realise.
The differences — source priority in \ac{sACN}, the mandatory data-loss timeout, and multicast subscription — are encapsulated entirely within the respective receiver classes and are invisible to the lighting engine.

\begin{table}[H]
  \caption{Comparison of Art-Net and \ac{sACN} protocol properties relevant to the receiver implementation.}
  \label{tab:proto_comparison}
  \centering
  \begin{tabularx}{\textwidth}{|l|X|X|}
    \hline
    \textbf{Property} & \textbf{Art-Net} & \textbf{sACN (ANSI E1.31)} \\ \hline \hline
    Transport          & UDP unicast / broadcast          & UDP unicast / multicast \\ \hline
    Port               & 6454                             & 5568 \\ \hline
    Universe range     & 0–32767                          & 1–63999 \\ \hline
    Channels/universe  & 512                              & 512 \\ \hline
    Source priority    & None                             & Per-source, 0–200 (default 100) \\ \hline
    Multi-source merge & Not specified                    & Single-winner, priority-based \\ \hline
    Data-loss handling & None                             & Mandatory \SI{2.5}{\second} timeout \\ \hline
    Subscription       & Unicast or broadcast             & Multicast group per universe \\ \hline
  \end{tabularx}
\end{table}

\subsubsection{Shared Data Layer}
\label{sec:net_shared_data}

The shared data layer consists of two components.
\texttt{universe\_config\_t} encodes the user-configurable mapping between network universes and \ac{LED} segments.
\texttt{dmx\_universe\_buf} is a double-buffered \ac{DMX} frame store that both protocol receivers write into.

\paragraph{Universe Configuration}
\texttt{universe\_config\_t} is written by Core~1 whenever the operator modifies the universe, start channel, or segment count, and is read by Core~0 (protocol receivers and lighting engine) as an atomic snapshot via the \texttt{portMUX\_TYPE} spinlock accessor described in Section~\ref{sec:shared_config}.
All boundary calculations that determine whether a segment's channel range overflows into a second universe are expressed as inline member functions, ensuring a single formula is applied consistently at every call site.

\begin{longlisting}
\begin{minted}{cpp}
static constexpr uint16_t DMX_ADDRESS_COUNT = 512; // valid: 1–512

struct universe_config_t {
  uint16_t start_universe; // first universe carrying segment data
  uint16_t start_channel;  // DMX start address (1–512)
  uint16_t num_segments;   // number of addressed LED segments

  // 3 RGB channels per segment + 1 global dimmer channel
  uint16_t used_channels() const { return num_segments * 3u + 1u; }

  bool spans_two_universes() const {
    return (start_channel + used_channels() - 1u) > DMX_ADDRESS_COUNT;
  }

  uint16_t end_universe() const {
    return start_universe + (spans_two_universes() ? 1u : 0u);
  }

  uint16_t end_channel() const {
    uint16_t last = start_channel + used_channels() - 1u;
    return spans_two_universes() ? last - DMX_ADDRESS_COUNT : last;
  }
};
\end{minted}
\caption{\texttt{universe\_config\_t}: universe mapping and boundary computation.}
\label{lst:universe_config}
\end{longlisting}

\paragraph{Double-Buffered Frame Store}
\texttt{dmx\_universe\_buf} holds one \ac{DMX} universe as a double buffer.
The write context — the Wi-Fi task that fires the network reception callback — always writes into the \emph{back buffer}.
The read context — the Core~0 output routine — always reads from the \emph{front buffer}.
After a complete frame has been written, the roles of the two buffers are exchanged by storing a new value into a single \texttt{std::atomic<uint8\_t>}.
This store maps to one native instruction and is therefore indivisible: the output routine either observes the old index or the new index, never a partial state.
The same byte encodes both the front-buffer index (bit~0) and a ready flag (bit~7), eliminating the need for a separate signalling variable.

\begin{figure}[H]
  \centering
  \begin{circuitikz}[>=latex]
    % Writer
    \draw[dashed, thick, fill=FrostBg, rounded corners] (0.0, 2.3) rectangle (4.8, 5.2);
    \node[anchor=north, font=\bfseries\small] at (2.4, 5.2) {Network Callback};
    \node[draw, fill=white, minimum width=3.8cm, minimum height=0.7cm, rounded corners, font=\small] (wrt)  at (2.4, 4.2) {Write to Back Buffer};
    \node[draw, fill=white, minimum width=3.8cm, minimum height=0.7cm, rounded corners, font=\small] (swap) at (2.4, 3) {Atomic Store (swap)};
    % Buffers
    \draw[dashed, thick, fill=FrostBg, rounded corners] (5.6, 0.3) rectangle (10.4, 5.2);
    \node[anchor=north, font=\bfseries\small] at (8.0, 5.2) {Double Buffer};
    \node[draw, fill=white, minimum width=3.8cm, minimum height=0.7cm, rounded corners, font=\small] (b0) at (8.0, 4.2) {Buffer 0 (512 B)};
    \node[draw, fill=white, minimum width=3.8cm, minimum height=0.7cm, rounded corners, font=\small] (b1) at (8.0, 3) {Buffer 1 (512 B)};
    \node[draw, fill=white, minimum width=3.8cm, minimum height=1.0cm, rounded corners, font=\small, align=center] (st) at (8.0, 1.5) {\texttt{atomic<uint8\_t>}\\state};
    % Reader
    \draw[dashed, thick, fill=FrostBg, rounded corners] (11.2, 2.3) rectangle (16.0, 5.2);
    \node[anchor=north, font=\bfseries\small] at (13.6, 5.2) {Output Routine};
    \node[draw, fill=white, minimum width=3.8cm, minimum height=0.7cm, rounded corners, font=\small] (rd)  at (13.6, 4.2) {Atomic Load (index)};
    \node[draw, fill=white, minimum width=3.8cm, minimum height=0.7cm, rounded corners, font=\small] (out) at (13.6, 3) {Drive LED String};
    % Arrows
    \draw[-latex, thick] (wrt.south)  -- (swap.north);
    \draw[-latex, thick] (rd.south)   -- (out.north);
    \draw[-latex, thick, rounded corners=4pt] (wrt.east)  -- (b0.west);
    \draw[-latex, thick, rounded corners=4pt] (wrt.east)  -- ++(0.8,0) |- (b1.west);
    \draw[-latex, thick, rounded corners=4pt] (swap.east) -- ++(0.6,0) |- (st.west);
    \draw[-latex, thick, rounded corners=4pt] (st.east)   -- ++(0.9,0) |- (rd.west);
    \draw[-latex, thick, rounded corners=4pt] (b0.east)   -- (rd.west);
    \draw[-latex, thick, rounded corners=4pt] (b1.east)   -- (out.west);
  \end{circuitikz}
  \caption{Double-buffered frame store operation.}
  \label{fig:dmx_double_buffer}
\end{figure}

As shown in Figure~\ref{fig:dmx_double_buffer}, the network callback writes \num{512} channel bytes into the back buffer and then atomically promotes it to the front by storing the new index and ready flag in a single instruction.
The output routine reads the atomic state to locate the front buffer and drives the \ac{LED} string from it without acquiring any lock.
A C++11 release-acquire pair on the atomic state establishes the necessary happens-before relationship.
\texttt{std::memory\_order\_release} in the callback guarantees that the completed \texttt{memcpy} is globally visible before the index swap is observed by any other execution context.
\texttt{std::memory\_order\_acquire} in the output routine guarantees that the buffer contents are read after the corresponding swap is observed.
On the Xtensa~LX7, both operations compile to a single native load or store instruction with no additional memory fence.

\begin{longlisting}
\begin{minted}{cpp}
class dmx_universe_buf {
public:
  static constexpr uint16_t CHANNEL_COUNT = DMX_ADDRESS_COUNT;

private:
  uint8_t frame_[2][CHANNEL_COUNT] = {};  // index 0 = DMX address 1
  std::atomic<uint8_t> state_{0}; // bit 0: front-buffer index; bit 7: ready flag

  static constexpr uint8_t READY_BIT = 0x80u;

public:
  // Called from the network reception callback.
  // Copies |count| channel bytes into the back buffer (DMX address 1 at index 0)
  // and atomically promotes the back buffer to the front.
  void write(const uint8_t* src, uint16_t count) {
    uint8_t back = (state_.load(std::memory_order_relaxed) & 1u) ^ 1u;
    memcpy(frame_[back], src,
      std::min(count, static_cast<uint16_t>(CHANNEL_COUNT)));
    state_.store(static_cast<uint8_t>(back | READY_BIT),
      std::memory_order_release);
  }

  // Called from the lighting engine.
  // Returns a pointer to the front buffer if a new frame is available,
  // or nullptr if no frame has arrived since the last call.
  // The channel byte for DMX address N is located at index N-1.
  const uint8_t* consume() {
    uint8_t s = state_.load(std::memory_order_acquire);
    if (!(s & READY_BIT)) return nullptr;
    state_.store(static_cast<uint8_t>(s & ~READY_BIT),
      std::memory_order_relaxed);
    return frame_[s & 1u];
  }
};
\end{minted}
\caption{\texttt{dmx\_universe\_buf}: double-buffered \ac{DMX} frame store with atomic index swap.}
\label{lst:dmx_universe_buf}
\end{longlisting}

The channel byte for \ac{DMX} address~$N$ is stored at array index~$N-1$, so address~1 maps to index~0 and address~512 maps to index~511.

\subsubsection{Network Manager}
\label{sec:net_manager}

The Network Manager is the sole component permitted to call \texttt{WiFi.begin()} or \texttt{WiFi.disconnect()}.
It implements a four-state \ac{FSM} and polls \texttt{WiFi.status()} once per Core~0 loop iteration.
State transitions post events to the Event Manager; protocol receivers react to \texttt{EVT\_WIFI\_CONNECTED} by calling \texttt{begin()} on their respective stacks, and to \texttt{EVT\_WIFI\_DISCONNECTED} by suspending packet processing.
On entry into \texttt{RECONNECTING}, a configurable back-off timer prevents rapid successive \texttt{WiFi.begin()} calls during a link failure.

\begin{figure}[H]
  \centering
  \begin{circuitikz}[>=latex]
    \node[draw, fill=white, minimum width=2.6cm, minimum height=0.8cm, rounded corners] (idle)   at (0.0,  1.5) {\texttt{IDLE}};
    \node[draw, fill=white, minimum width=2.8cm, minimum height=0.8cm, rounded corners] (conn)   at (5.0,  1.5) {\texttt{CONNECTING}};
    \node[draw, fill=FrostBg, minimum width=2.8cm, minimum height=0.8cm, rounded corners] (up)     at (10.5, 1.5) {\texttt{CONNECTED}};
    \node[draw, fill=white, minimum width=2.8cm, minimum height=0.8cm, rounded corners] (reconn) at (10.5, 0.0) {\texttt{RECONNECTING}};

    \draw[-latex, thick] (idle.east) -- (conn.west)
      node[midway, above, font=\scriptsize] {start()};
    \draw[-latex, thick] (conn.east) -- (up.west)
      node[midway, above, font=\scriptsize] {WL\_CONNECTED};
    \draw[-latex, thick] (up.south) -- (reconn.north)
      node[midway, right, font=\scriptsize] {link lost};
    \draw[-latex, thick, rounded corners=4pt] (reconn.west) -- ++(-1.5,0) |- (up.west);
    \node[font=\scriptsize, right] at (7.5, 0.75) {WL\_CONNECTED};
    \draw[-latex, thick, rounded corners=4pt] (up.north) -- ++(0,0.6) -| (idle.north);
    \node[font=\scriptsize, above] at (5.25, 2.5) {stop()};
  \end{circuitikz}
  \caption{Network Manager \ac{FSM}.}
  \label{fig:wifi_fsm}
\end{figure}

\texttt{stop()} is valid from any active state; the arc from \texttt{CONNECTED} is shown for clarity.
On repeated \texttt{RECONNECTING} cycles, \texttt{WiFi.begin()} is re-issued after the back-off delay elapses.

\begin{longlisting}
\begin{minted}{cpp}
enum class wifi_state_t { IDLE, CONNECTING, CONNECTED, RECONNECTING };

class network_manager {
private:
  wifi_state_t state_        = wifi_state_t::IDLE;
  uint32_t     reconnect_at_ = 0;
  const char*  ssid_;
  const char*  password_;

  static constexpr uint32_t RECONNECT_DELAY_MS = 5000u;

public:
  network_manager(const char* ssid, const char* password)
    : ssid_(ssid), password_(password) {}

  void start() {
    WiFi.begin(ssid_, password_);
    state_ = wifi_state_t::CONNECTING;
  }

  void stop() {
    WiFi.disconnect();
    WiFi.mode(WIFI_OFF);
    state_ = wifi_state_t::IDLE;
    event_manager.post(EVT_WIFI_STOPPED);
  }

  // Polled once per Core 0 main-loop iteration.
  void poll() {
    const bool linked = (WiFi.status() == WL_CONNECTED);
    switch (state_) {
      case wifi_state_t::CONNECTING:
        if (linked) {
          state_ = wifi_state_t::CONNECTED;
          event_manager.post(EVT_WIFI_CONNECTED);
        }
        break;
      case wifi_state_t::CONNECTED:
        if (!linked) {
          state_        = wifi_state_t::RECONNECTING;
          reconnect_at_ = millis() + RECONNECT_DELAY_MS;
          event_manager.post(EVT_WIFI_DISCONNECTED);
        }
        break;
      case wifi_state_t::RECONNECTING:
        if (linked) {
          state_ = wifi_state_t::CONNECTED;
          event_manager.post(EVT_WIFI_CONNECTED);
        }
        else if (millis() >= reconnect_at_) {
          WiFi.begin(ssid_, password_);
          reconnect_at_ = millis() + RECONNECT_DELAY_MS;
        }
        break;
      default: break;
    }
  }
};
\end{minted}
\caption{Network Manager implementation.}
\label{lst:network_manager}
\end{longlisting}

Connection recovery is handled autonomously: after \texttt{EVT\_WIFI\_DISCONNECTED} is posted, the Manager re-issues \texttt{WiFi.begin()} every \SI{5}{\second} until the link is restored, then posts \texttt{EVT\_WIFI\_CONNECTED} so the active protocol receiver can resume.

\subsubsection{Art-Net Receiver}
\label{sec:artnet_receiver}

The Art-Net receiver wraps the ArtnetWifi library and is responsible for universe filtering and frame dispatch.
On initialisation it registers a lambda with the library as the ArtDMX callback; the library invokes this lambda whenever a valid ArtDMX packet arrives on \ac{UDP} port~6454.
The callback reads a spinlock snapshot of \texttt{universe\_config\_t}, compares the incoming universe number against the configured start and end universes, and writes the payload into the appropriate frame buffer.
The receiver holds no protocol state of its own; all Art-Net framing and sequencing is handled by the library.

\begin{longlisting}
\begin{minted}{cpp}
class artnet_receiver {
private:
  ArtnetWifi        artnet_;
  dmx_universe_buf* primary_;
  dmx_universe_buf* secondary_;

public:
  void init(dmx_universe_buf* primary, dmx_universe_buf* secondary) {
    primary_   = primary;
    secondary_ = secondary;
    artnet_.setArtDmxCallback(
      [this](uint16_t universe, uint16_t length,
             uint8_t /*seq*/, uint8_t* data) {
        on_dmx_frame(universe, length, data);
      });
  }

  void begin() { artnet_.begin(); } // called on EVT_WIFI_CONNECTED
  void poll()  { artnet_.read();  } // called once per Core 0 loop iteration

private:
  void on_dmx_frame(uint16_t universe, uint16_t length, const uint8_t* data) {
    // Snapshot avoids holding the spinlock across the memcpy inside write().
    const universe_config_t cfg = config_get();

    if (universe == cfg.start_universe) {
      primary_->write(data, length);
    } else if (cfg.spans_two_universes() && universe == cfg.end_universe()) {
      secondary_->write(data, length);
    }
  }
};
\end{minted}
\caption{Art-Net receiver implementation.}
\label{lst:artnet_receiver}
\end{longlisting}

The lambda captures \texttt{this} so the stateless closure has access to the buffer pointers.
Taking a spinlock snapshot before calling \texttt{write()} confines the critical section to a struct copy, keeping the spinlock hold time independent of the \num{512}-byte \texttt{memcpy}.

\subsubsection{sACN Receiver}
\label{sec:sacn_receiver}

The \ac{sACN} receiver implements ANSI~E1.31 reception and adds three capabilities beyond the Art-Net receiver: multicast subscription, single-source selection with priority-based takeover, and mandatory data-loss detection.

\paragraph{Multicast Subscription}
\ac{sACN} transmitters address each universe through an \ac{IP} multicast group.
The multicast address for universe~$u$ is \texttt{239.255.(u >> 8).(u \& 0xFF)}.
On \texttt{begin()}, the receiver joins the multicast group for \texttt{start\_universe} and, when the channel range spans two universes, the group for \texttt{end\_universe} as well.
Group membership is refreshed whenever the user changes the universe configuration.

\paragraph{Source Selection and Priority}
\ac{sACN} permits multiple transmitters to address the same universe simultaneously, each carrying a priority value in the range 0–200 (default 100).
The receiver maintains one \emph{active source} identified by its 16-byte component identifier.
Packets from the active source are always accepted; packets from a different source are accepted only if they carry a strictly higher priority, at which point that source becomes the new active source.
This single-winner policy avoids per-slot priority merging while meeting the baseline interoperability requirement of ANSI~E1.31.

\paragraph{Data-Loss Detection}
ANSI~E1.31 §6.7.1 requires that a receiver treat the absence of valid packets for \SI{2.5}{\second} as a data-loss condition.
The timestamp of each accepted packet is recorded; the elapsed time is evaluated on every poll call.
On timeout, the active source record is cleared and \texttt{EVT\_SACN\_DATA\_LOSS} is posted to the Event Manager.
The Core~1 event handler responds with a hold-last-look policy, preserving the last rendered frame and preventing an abrupt blackout.

\begin{longlisting}
\begin{minted}{cpp}
struct sacn_source_t {
  uint8_t  cid[16];           // 16-byte component identifier (UUID-style)
  uint8_t  priority;
  uint32_t last_pkt_ms;
};

class sacn_receiver {
private:
  WiFiUDP udp_;
  dmx_universe_buf* primary_;
  dmx_universe_buf* secondary_;
  sacn_source_t source_     = {};
  bool has_source_ = false;

  static constexpr uint16_t E131_PORT = 5568u;
  static constexpr uint32_t DATA_LOSS_MS = 2500u;
  static constexpr uint16_t MIN_PKT_LEN = 127u; // 126-byte header + 1 data byte

public:
  void begin(dmx_universe_buf* primary, dmx_universe_buf* secondary) {
    primary_   = primary;
    secondary_ = secondary;
    udp_.begin(E131_PORT);
    const universe_config_t cfg = config_get();
    join_multicast(cfg.start_universe);
    if (cfg.spans_two_universes()) join_multicast(cfg.end_universe());
  }

  void stop() { udp_.stop(); has_source_ = false; }

  void poll() {
    check_data_loss();
    const int len = udp_.parsePacket();
    if (len < MIN_PKT_LEN) return;
    uint8_t pkt[638];
    if (udp_.read(pkt, sizeof(pkt)) < MIN_PKT_LEN) return;
    if (!valid_preamble(pkt)) return;

    const uint16_t universe = (uint16_t(pkt[113]) << 8) | pkt[114];
    const uint8_t  priority = pkt[108];
    const uint16_t slots = (uint16_t(pkt[123]) << 8) | pkt[124];
    const uint8_t* dmx = &pkt[126];  // first channel value; start code excluded

    if (slots < 2u || !accept_source(&pkt[22], priority)) return;

    const universe_config_t cfg = config_get();
    if (universe == cfg.start_universe)
      primary_->write(dmx, slots - 1u);
    else if (cfg.spans_two_universes() && universe == cfg.end_universe())
      secondary_->write(dmx, slots - 1u);
  }

private:
  void join_multicast(uint16_t u) {
    udp_.beginMulticast(IPAddress(239u, 255u, u >> 8, u & 0xFFu), E131_PORT);
  }

  bool accept_source(const uint8_t* cid, uint8_t priority) {
    const uint32_t now = millis();
    if (!has_source_) {
      memcpy(source_.cid, cid, 16);
      source_.priority = priority; source_.last_pkt_ms = now;
      return has_source_ = true;
    }
    if (memcmp(source_.cid, cid, 16) == 0) {
      source_.priority = priority; source_.last_pkt_ms = now;
      return true;
    }
    if (priority > source_.priority) {
      memcpy(source_.cid, cid, 16);
      source_.priority = priority; source_.last_pkt_ms = now;
      return true;
    }
    return false;
  }

  void check_data_loss() {
    if (has_source_ && (millis() - source_.last_pkt_ms > DATA_LOSS_MS)) {
      has_source_ = false;
      event_manager.post(EVT_SACN_DATA_LOSS);
    }
  }

  static bool valid_preamble(const uint8_t* p) {
    if (p[0] != 0x00u || p[1] != 0x10u) return false;
    static const uint8_t ACN_ID[12] =
      {0x41,0x53,0x43,0x2D,0x45,0x31,0x2E,0x31,0x37,0x00,0x00,0x00};
    return memcmp(&p[4], ACN_ID, 12) == 0;
  }
};
\end{minted}
\caption{\ac{sACN} receiver class.}
\label{lst:sacn_receiver}
\end{longlisting}

\texttt{accept\_source()} implements the single-winner selection policy: it always accepts the known active source and replaces it only when a competing source arrives with strictly higher priority.
\texttt{valid\_preamble()} checks the ANSI~E1.31 preamble-size field and the ACN packet identifier, rejecting malformed or non-E1.31 datagrams before any further parsing occurs.

\subsubsection{Lighting Engine Integration}
\label{sec:net_lighting_integration}

The lighting engine interacts with the shared data layer exclusively through two operations: \texttt{config\_get()} for a spinlock-protected snapshot of \texttt{universe\_config\_t}, and \texttt{dmx\_universe\_buf::consume()} on each frame buffer to retrieve the latest channel data.
\texttt{consume()} returns \texttt{nullptr} when no new frame has arrived since the last call, allowing the engine to skip rendering entirely and avoid redundant \ac{LED} updates.
When a channel range spans two universes, the engine waits until both the primary and secondary buffers have delivered a new frame before rendering, ensuring that an \ac{LED} string is never driven with a split-universe mix of an old frame and a new one.
The function is entirely protocol-agnostic: the data in the buffers may have originated from an Art-Net or an \ac{sACN} packet; the engine cannot and need not distinguish between them.

\begin{longlisting}
\begin{minted}{cpp}
void lighting_engine_tick(dmx_universe_buf& primary, dmx_universe_buf& secondary) {
  const universe_config_t cfg = config_get();

  const uint8_t* pdata = primary.consume();
  if (!pdata) return;  // no new primary-universe frame

  const uint8_t* sdata = nullptr;
  if (cfg.spans_two_universes()) {
    sdata = secondary.consume();
    if (!sdata) return;  // hold until the secondary universe also arrives
  }

  const uint16_t pixels_per_seg  = NUM_PIXEL / cfg.num_segments;
  uint16_t       channel_idx     = cfg.start_channel - 1u;  // DMX address → 0-based
  uint16_t       led_idx         = 0;

  for (uint16_t seg = 0; seg < cfg.num_segments; ++seg) {
    // Resolve each channel index against the correct universe buffer.
    // Indices 0–(DMX_ADDRESS_COUNT-1) map to pdata; overflow maps to sdata.
    auto ch = [&](uint16_t idx) -> uint8_t {
      return (idx < DMX_ADDRESS_COUNT) ? pdata[idx] : sdata[idx - DMX_ADDRESS_COUNT];
    };

    const uint8_t r = ch(channel_idx++);
    const uint8_t g = ch(channel_idx++);
    const uint8_t b = ch(channel_idx++);

    // Global dimmer occupies the last channel in the configured range.
    const uint8_t dim = ch(cfg.end_channel() - 1u + (cfg.start_channel - 1u));

    CRGB color(r, g, b);
    color.nscale8_video(dim);
    for (uint16_t px = 0; px < pixels_per_seg; ++px)
      pixels.setPixelColor(led_idx++, pixels.Color(color.r, color.g, color.b));
  }
  pixels.show();
}
\end{minted}
\caption{Lighting engine tick function.}
\label{lst:lighting_engine_tick}
\end{longlisting}

The lambda \texttt{ch()} resolves a 0-based channel index against the primary or secondary buffer, transparently handling the universe boundary.
The function is protocol-agnostic: it does not inspect the origin of the data in the frame buffers.
